% 第 4 题解析：P 点电场强度叠加示意
% +q 在 (a,0)、-q 在 (0,a)：P 点合场强 E1 方向左上 45°；Q 在 (0,2a)，
% 其场强 EQ 方向右下 45°（背离 Q），与 E1 等大反向使 P 点场强为零。
\begin{tikzpicture}[font=\small, >=Stealth, scale=1.0]
  \def\a{1.5}
  % 坐标轴
  \draw[->] (-0.8,0) -- (3.4,0) node[right] {$x$};
  \draw[->] (0,-0.8) -- (0,3.6) node[above] {$y$};
  % 电荷位置
  \fill ({\a},0) circle (2pt) node[below] {$(a,0)\ +q$};
  \fill (0,{\a}) circle (2pt) node[left] {$(0,a)\ -q$};
  \fill ({\a},{\a}) circle (2pt) node[above right] {$P(a,a)$};
  \fill (0,{2*\a}) circle (2pt) node[above left] {$Q$};
  % 场强矢量：E1（左上 45°，由 -q 与 +q 场强合成）与 EQ（右下 45°）
  \draw[->, very thick, blue] ({\a},{\a}) -- ({-0.45*\a},{1.45*\a})
      node[midway, above left] {$\gv{E}_{1}$};
  \draw[->, very thick, red] ({\a},{\a}) -- ({1.45*\a},{0.55*\a})
      node[midway, below right] {$\gv{E}_{Q}$};
  % 距离示意
  \draw[dashed] (0,{\a}) -- ({\a},{\a});
  \draw[dashed] ({\a},0) -- ({\a},{\a});
  \draw[dashed] ({\a},{\a}) -- (0,{2*\a});
  \node[below left] at (0.3,0.2) {$O$};
\end{tikzpicture}
